%% sol-ch06.tex — Solutions for Chapter 6: Regular Expressions

\chapter{Regular Expressions --- Solutions}

\begin{enumerate}[{6.}1.]

%% 6.1
\item Even-length words over $\{\pmb{a},\pmb{b}\}$:
\[
  ((\pmb{a}\cup\pmb{b})(\pmb{a}\cup\pmb{b}))^*
\]

%% 6.2
\item Words of length $\ge 2$ over $\{\pmb{a},\pmb{b}\}$:
\[
  (\pmb{a}\cup\pmb{b})(\pmb{a}\cup\pmb{b})(\pmb{a}\cup\pmb{b})^*
\]

%% 6.3
\item Words with $|x|_{\pmb{a}}=|x|_{\pmb{b}}$ over $\{\pmb{a},\pmb{b}\}$: This language is \emph{not} regular (not FAD).  By the pumping lemma, if it were FAD with $n$ states, take $w=\pmb{a}^n\pmb{b}^n$; any pump within the first $n$ characters (all $\pmb{a}$s) would unbalance the counts.  Hence no regular expression exists.

%% 6.4
\item Odd-length words over $\{\pmb{a},\pmb{b},\pmb{c}\}$ not ending in $\pmb{b}$:

A word of odd length not ending in $\pmb{b}$ has an even-length prefix (possibly empty) followed by a non-$\pmb{b}$ letter:
\[
  ((\pmb{a}\cup\pmb{b}\cup\pmb{c})(\pmb{a}\cup\pmb{b}\cup\pmb{c}))^*(\pmb{a}\cup\pmb{c})
\]
(When the Kleene star produces $k$ pairs, the total length is $2k+1$, which is odd.)

%% 6.5
\item Words over $\{\pmb{a},\pmb{b},\pmb{c}\}$ not containing $\pmb{cc}$:
Words with no two consecutive $\pmb{c}$s are those where every $\pmb{c}$ is followed by a non-$\pmb{c}$ (or is at the end).  Let $N=\pmb{a}\cup\pmb{b}$ (non-$\pmb{c}$ letters).
\[
  (N \cup \pmb{c}N)^*(N \cup \pmb{c} \cup \epsilon)
\]

%% 6.6
\item Words over $\{\pmb{a},\pmb{b},\pmb{c}\}$ containing $\pmb{cc}$:
\[
  (\pmb{a}\cup\pmb{b}\cup\pmb{c})^*\pmb{cc}(\pmb{a}\cup\pmb{b}\cup\pmb{c})^*
\]

%% 6.7
\item Words over $\{\pmb{a},\pmb{b},\pmb{c}\}$ with no $\pmb{c}$s:
\[
  (\pmb{a}\cup\pmb{b})^*
\]

%% 6.8
\item Let $\Sigma=\{\pmb{0},\pmb{1}\}$.

\begin{enumerate}[(a)]
  \item $L_1=\{x\mid|x|\bmod 3=2\}$: words whose length is congruent to 2 modulo 3:
  \[
    (\pmb{0}\cup\pmb{1})(\pmb{0}\cup\pmb{1})((\pmb{0}\cup\pmb{1})(\pmb{0}\cup\pmb{1})(\pmb{0}\cup\pmb{1}))^*
  \]

  \item $L_2=\Sigma^*\setminus\{w\mid w\text{ ends in }\pmb{1}\}$: words that do \emph{not} end in $\pmb{1}$, i.e., words that end in $\pmb{0}$ or equal $\lambda$:
  \[
    \epsilon\cup(\pmb{0}\cup\pmb{1})^*\pmb{0}
  \]

  \item $L_3=\{y\mid|y|_{\pmb{0}}>|y|_{\pmb{1}}\}$: this language is \emph{not}
  regular.  If it were regular, then its intersection with the regular language
  $\pmb{0}^*\pmb{1}^*$ would also be regular.  But
  \[
    L_3\cap\pmb{0}^*\pmb{1}^*=\{\pmb{0}^m\pmb{1}^n\mid m>n\},
  \]
  and this language is not regular by the pumping lemma: for pumping length
  $p$, the word $\pmb{0}^{p+1}\pmb{1}^p$ belongs to it, but pumping down any
  nonempty part from the initial $\pmb{0}$-block yields a word with at most as
  many $\pmb{0}$s as $\pmb{1}$s.  Hence no regular expression exists.
\end{enumerate}

%% 6.9
\item Let $\Sigma=\{\pmb{a},\pmb{b},\pmb{c}\}$.  Let $\Sigma^*=(\pmb{a}\cup\pmb{b}\cup\pmb{c})^*$.

\begin{enumerate}[(a)]
  \item $L_1=\{x\mid|x|_{\pmb{a}}\text{ odd}\wedge|x|_{\pmb{b}}\text{ even}\}$:

  Let $C=\pmb{c}^*$, $A=\pmb{a}C$, and $B=\pmb{b}C$.  Then the non-$\pmb{c}$
  letters must contribute odd $\pmb{a}$-parity and even $\pmb{b}$-parity.  One
  exact regular expression is:
  \[
    C\bigl(AA\cup BB\cup(AB\cup BA)(AA\cup BB)^*(AB\cup BA)\bigr)^*
    \bigl(A\cup(AB\cup BA)(AA\cup BB)^*B\bigr)
  \]

  \item $L_2=\{y\mid|y|_{\pmb{c}}\text{ even}\vee|y|_{\pmb{b}}\text{ odd}\}$:
  \[
    \bigl((\pmb{a}\cup\pmb{b})^*\pmb{c}(\pmb{a}\cup\pmb{b})^*\pmb{c}\bigr)^*(\pmb{a}\cup\pmb{b})^*
    \cup
    \bigl((\pmb{a}\cup\pmb{c})^*\pmb{b}(\pmb{a}\cup\pmb{c})^*\pmb{b}\bigr)^*(\pmb{a}\cup\pmb{c})^*\pmb{b}(\pmb{a}\cup\pmb{c})^*
  \]
  The first expression gives all words with an even number of $\pmb{c}$s; the
  second gives all words with an odd number of $\pmb{b}$s.

  \item $L_3=\{z\mid|z|_{\pmb{a}}\text{ even}\}$:
  \[
    ((\pmb{b}\cup\pmb{c})^*\pmb{a}(\pmb{b}\cup\pmb{c})^*\pmb{a})^*(\pmb{b}\cup\pmb{c})^*
  \]

  \item $L_4=\{z\mid|z|_{\pmb{c}}\text{ prime}\}$: This language is not regular.
  If it were FAD, let $p$ be a pumping length and choose a prime $q>p$.
  Then $\pmb{c}^q\in L_4$.  Write
  \[
    \pmb{c}^q=uvw
  \]
  with $|uv|\le p$ and $|v|\ge 1$.  Since the word is unary, let $|v|=m\ge 1$.
  Pump with $i=q+1$.  The pumped word has length
  \[
    q+qm=q(1+m),
  \]
  which is composite.  Hence $uv^{q+1}w\notin L_4$, contradicting the pumping
  lemma.  Therefore $L_4$ has no regular expression.

  \item $L_5=\{x\mid\pmb{abc}\text{ is a substring}\}$:
  \[
    (\pmb{a}\cup\pmb{b}\cup\pmb{c})^*\pmb{abc}(\pmb{a}\cup\pmb{b}\cup\pmb{c})^*
  \]

  \item $L_6=\{x\mid\pmb{acaba}\text{ is a substring}\}$:
  \[
    (\pmb{a}\cup\pmb{b}\cup\pmb{c})^*\pmb{acaba}(\pmb{a}\cup\pmb{b}\cup\pmb{c})^*
  \]

  \item $L_7=\{x\mid|x|_{\pmb{a}}\equiv 0\pmod{3}\}$:
  \[
    (\pmb{b}\cup\pmb{c})^*(\pmb{a}(\pmb{b}\cup\pmb{c})^*\pmb{a}(\pmb{b}\cup\pmb{c})^*\pmb{a}(\pmb{b}\cup\pmb{c})^*)^*
  \]
\end{enumerate}

%% 6.10
\item $\Psi=\{x\mid x\text{ begins with }\pmb{d}\vee x\text{ contains }\pmb{bb}\}$ over $\{\pmb{a},\pmb{b},\pmb{d}\}$:
\[
  \pmb{d}(\pmb{a}\cup\pmb{b}\cup\pmb{d})^*
  \cup
  (\pmb{a}\cup\pmb{b}\cup\pmb{d})^*\pmb{bb}(\pmb{a}\cup\pmb{b}\cup\pmb{d})^*
\]

%% 6.11
\item $\Phi=\{x\mid\text{every }\pmb{b}\text{ in }x\text{ is immediately followed by }\pmb{c}\}$ over $\{\pmb{a},\pmb{b},\pmb{c}\}$:
Words with no isolated $\pmb{b}$ (every $\pmb{b}$ must be followed immediately by $\pmb{c}$).  Let $S=\pmb{a}\cup\pmb{c}\cup\pmb{bc}$ (safe building blocks):
\[
  (\pmb{a}\cup\pmb{c}\cup\pmb{bc})^*
\]

%% 6.12
\item Numbers divisible by 3 over $\{0,\ldots,9\}$:

A number is divisible by 3 iff the sum of its digits is divisible by 3.  Let
\[
D_0=(\pmb{0}\cup\pmb{3}\cup\pmb{6}\cup\pmb{9}),\quad
D_1=(\pmb{1}\cup\pmb{4}\cup\pmb{7}),\quad
D_2=(\pmb{2}\cup\pmb{5}\cup\pmb{8}).
\]
Using the $3$-state DFA that tracks the running digit sum modulo $3$, the
terminal-set equations are
\begin{align*}
  X_0 &= \epsilon \cup D_0X_0 \cup D_1X_1 \cup D_2X_2 \\
  X_1 &= D_2X_0 \cup D_0X_1 \cup D_1X_2 \\
  X_2 &= D_1X_0 \cup D_2X_1 \cup D_0X_2 .
\end{align*}
Applying Theorems~6.2 and~6.1 yields the following exact regular expression for
the language:
\[
\begin{aligned}
X_0={}&((D_0)^*\cup(D_0)^*D_1((D_0\cup D_2(D_0)^*D_1))^*D_2(D_0)^*) \\
&{}\cup((D_0)^*D_2\cup(D_0)^*D_1((D_0\cup D_2(D_0)^*D_1))^*(D_1\cup D_2(D_0)^*D_2)) \\
&\qquad\cdot(((D_0\cup D_1(D_0)^*D_2)\cup(D_2\cup D_1(D_0)^*D_1)((D_0\cup D_2(D_0)^*D_1))^*(D_1\cup D_2(D_0)^*D_2)))^* \\
&\qquad\cdot(D_1(D_0)^*\cup(D_2\cup D_1(D_0)^*D_1)((D_0\cup D_2(D_0)^*D_1))^*D_2(D_0)^*)) .
\end{aligned}
\]

%% 6.13
\item $K=\{x\mid x\text{ is divisible by }5\}$ over decimal digits:

A decimal number is divisible by 5 iff it ends in $\pmb{0}$ or $\pmb{5}$.
\[
  K = (\pmb{0}\cup\pmb{1}\cup\cdots\cup\pmb{9})^*(\pmb{0}\cup\pmb{5})
\]

%% 6.14
\item \textbf{Build NDFA for $\pmb{b}\cup\pmb{a}^*\pmb{c}$ using exact constructs of Chapter~5.}

Use the same state names and unsimplified constructions as in
Examples~6.4--6.6.

\medskip
\noindent\textbf{Step 1: Build $\{\pmb{a}\}$ and apply Theorem~5.5.}
Let
\[
A_{\pmb{a}}=\langle\Sigma,\{t_0,t_1\},\{t_0\},\delta_{\pmb{a}},\{t_1\}\rangle
\]
with $\delta_{\pmb{a}}(t_0,\pmb{a})=\{t_1\}$ and all other transitions empty.
Then Theorem~5.5 gives an NDFA for $\pmb{a}^*$:
\[
A_{\pmb{a}^*}=\langle\Sigma,\{t_0,t_1,q_0\},\{t_0,q_0\},\delta_*,\{t_1,q_0\}\rangle,
\]
where
\[
\delta_*(t_0,\pmb{a})=\{t_1\},\qquad
\delta_*(t_1,\pmb{a})=\{t_1\},
\]
and every other transition is empty.

\medskip
\noindent\textbf{Step 2: Concatenate with $\{\pmb{c}\}$ using the pattern of Example~6.6.}
Let
\[
A_{\pmb{c}}=\langle\Sigma,\{r_0,r_1\},\{r_0\},\delta_{\pmb{c}},\{r_1\}\rangle
\]
with $\delta_{\pmb{c}}(r_0,\pmb{c})=\{r_1\}$ and all other transitions empty.
Using the exact concatenation construction, form
\[
A_{\pmb{a}^*\pmb{c}}
=\langle\Sigma,\{t_0,t_1,q_0,r_0,r_1\},\{t_0,q_0\},\delta_\bullet,\{r_1\}\rangle,
\]
where $\delta_\bullet$ agrees with the old transitions and, because $t_1$ and
$q_0$ are final in $A_{\pmb{a}^*}$, they also mimic the transitions out of
$r_0$.  Thus
\[
\delta_\bullet(t_0,\pmb{a})=\{t_1\},\qquad
\delta_\bullet(t_1,\pmb{a})=\{t_1\},\qquad
\delta_\bullet(t_1,\pmb{c})=\{r_1\},
\]
\[
\delta_\bullet(q_0,\pmb{c})=\{r_1\},\qquad
\delta_\bullet(r_0,\pmb{c})=\{r_1\},
\]
and every other transition is empty.  As in Example~6.6, the state $r_0$
remains disconnected.

\medskip
\noindent\textbf{Step 3: Union with $\{\pmb{b}\}$ using Theorem~5.2.}
Let
\[
A_{\pmb{b}}=\langle\Sigma,\{s_0,s_1\},\{s_0\},\delta_{\pmb{b}},\{s_1\}\rangle
\]
with $\delta_{\pmb{b}}(s_0,\pmb{b})=\{s_1\}$ and all other transitions empty.
Then Theorem~5.2 yields
\[
A=
\langle
\Sigma,
\{t_0,t_1,q_0,r_0,r_1,s_0,s_1\},
\{t_0,q_0,s_0\},
\delta,
\{r_1,s_1\}
\rangle,
\]
where $\delta$ agrees with $\delta_\bullet$ on
$\{t_0,t_1,q_0,r_0,r_1\}$, agrees with $\delta_{\pmb{b}}$ on
$\{s_0,s_1\}$, and has no other transitions.  This is the required unsimplified
NDFA for $\pmb{b}\cup\pmb{a}^*\pmb{c}$.

%% 6.15
\item \textbf{Inequalities in Lemma~6.1:}

\begin{enumerate}[(a)]
  \item $R_1\cup\epsilon\neq R_1$: Let $R_1=\pmb{a}$.  Then $R_1\cup\epsilon=\{\lambda,\pmb{a}\}\neq\{\pmb{a}\}=R_1$.
  \item $R_1\cdot R_2\neq R_2\cdot R_1$: Let $R_1=\pmb{a}$, $R_2=\pmb{b}$.  $\pmb{ab}\neq\pmb{ba}$.
  \item $R_1\cdot R_1\neq R_1$: Let $R_1=\pmb{a}$.  $\pmb{aa}\neq\pmb{a}$.
  \item $R_1\cup(R_2\cdot R_3)\neq(R_1\cup R_2)\cdot(R_1\cup R_3)$: Let $R_1=\pmb{a}$, $R_2=\pmb{b}$, $R_3=\pmb{c}$.  Then
  \[
    R_1\cup(R_2\cdot R_3)=\{\pmb{a},\pmb{bc}\},
    \qquad
    (R_1\cup R_2)\cdot(R_1\cup R_3)=\{\pmb{a},\pmb{b}\}\cdot\{\pmb{a},\pmb{c}\}
    =\{\pmb{aa},\pmb{ac},\pmb{ba},\pmb{bc}\}.
  \]
  These two sets are not equal.
  \item $(R_1\cdot R_2)^*\neq(R_1^*\cdot R_2^*)^*$: Let $R_1=\pmb{a}$, $R_2=\pmb{b}$.  $(ab)^*=\{\lambda,\pmb{ab},\pmb{abab},\ldots\}$.  $(a^*b^*)^*=(\pmb{a}\cup\pmb{b})^*$.  Not equal.
  \item $(R_1\cdot R_2)^*\neq(R_1^*\cup R_2^*)^*$: Same example.  $(R_1^*\cup R_2^*)^*=(a^*\cup b^*)^*=(a\cup b)^*$.  Not equal to $(\pmb{ab})^*$.
\end{enumerate}

\textbf{Examples where equality holds:}
\begin{enumerate}[(a)]
  \setcounter{enumii}{6}
  \item $R_1\cup\epsilon=R_1$: let $R_1=\pmb{a}^*$.  Then $\pmb{a}^*\cup\epsilon=\pmb{a}^*$.
  \item $R_1\cdot R_2=R_2\cdot R_1$: let $R_1=\pmb{a}$ and $R_2=\pmb{aa}$.  Then $R_1\neq R_2$, neither is $\epsilon$, and both products equal $\pmb{aaa}$.
  \item $R_1\cdot R_1=R_1$: let $R_1=\Sigma^*$.  Then $\Sigma^*\Sigma^*=\Sigma^*$.
  \item $R_1\cup(R_2\cdot R_3)=(R_1\cup R_2)\cdot(R_1\cup R_3)$: let $R_1=(\pmb{a}\cup\pmb{b})^*$, $R_2=\pmb{a}$, and $R_3=\pmb{b}$.  These three sets are distinct, and both sides equal $(\pmb{a}\cup\pmb{b})^*$.
  \item $(R_1\cdot R_2)^*=(R_1^*\cdot R_2^*)^*$: let $R_1=\pmb{a}^*$ and $R_2=\pmb{a}\pmb{a}^*$.  Then $R_1\neq R_2$, neither is $\epsilon$, and both sides equal $\pmb{a}^*$.
  \item $(R_1\cdot R_2)^*=(R_1^*\cup R_2^*)^*$: with the same choice $R_1=\pmb{a}^*$ and $R_2=\pmb{a}\pmb{a}^*$, both sides again equal $\pmb{a}^*$.
\end{enumerate}

%% 6.16
\item \textbf{Prove the equalities of Lemma~6.1.}

Let $L_i=L(R_i)$.  Each identity follows directly from the definitions.

\begin{enumerate}[(a)]
  \item $R_1\cup\emptyset=R_1$: $L_1\cup\emptyset=L_1$.
  \item $R_1\cdot\epsilon=R_1=\epsilon\cdot R_1$: concatenating with $\{\lambda\}$ does not change a word.
  \item $R_1\cdot\emptyset=\emptyset=\emptyset\cdot R_1$: there is no second factor to form a concatenation.
  \item $R_1\cup R_2=R_2\cup R_1$: union is commutative.
  \item $R_1\cup R_1=R_1$: union is idempotent.
  \item $R_1\cup(R_2\cup R_3)=(R_1\cup R_2)\cup R_3$: union is associative.
  \item $R_1\cdot(R_2\cdot R_3)=(R_1\cdot R_2)\cdot R_3$: both sides are all words of the form $xyz$ with $x\in L_1$, $y\in L_2$, $z\in L_3$.
  \item $R_1\cdot(R_2\cup R_3)=(R_1\cdot R_2)\cup(R_1\cdot R_3)$: a word on the left is $xy$ with $x\in L_1$ and $y\in L_2\cup L_3$, so $y\in L_2$ or $y\in L_3$.
  \item $\epsilon^*=\epsilon$: concatenating any number of empty words still gives $\lambda$.
  \item $\emptyset^*=\epsilon$: by definition, $L^*=\bigcup_{n\ge 0}L^n$, and $\emptyset^0=\epsilon$ while $\emptyset^n=\emptyset$ for $n\ge 1$.
  \item $(R_1\cup R_2)^*=(R_1^*\cup R_2^*)^*$: every factor from $L_1\cup L_2$ already lies in $L_1^*\cup L_2^*$, so the left side is contained in the right; conversely, each factor from $L_1^*$ or $L_2^*$ is itself a finite concatenation of words from $L_1\cup L_2$, so the right side is contained in the left.
  \item $(R_1\cup R_2)^*=(R_1^*\cdot R_2^*)^*$: a block from $L_1^*L_2^*$ is a finite string of words from $L_1$ followed by a finite string of words from $L_2$, hence belongs to $(L_1\cup L_2)^*$; the reverse inclusion holds because each factor from $L_1\cup L_2$ can be viewed as a block with either the $L_1^*$ or $L_2^*$ part empty.
  \item $(R_1^*)^*=R_1^*$: $L_1^*$ is already closed under finite concatenation.
  \item $R_1^*\cdot R_1^*=R_1^*$: concatenating two finite concatenations of words from $L_1$ still yields a finite concatenation of words from $L_1$, and every word of $L_1^*$ is obtained by taking the second factor to be $\lambda$.
\end{enumerate}

%% 6.17
\item \begin{enumerate}[(a)]
  \item \textbf{Non-uniqueness of $A^*E$ when $\lambda\in A$.}

  Let $A=\{\lambda,\pmb{a}\}=\epsilon\cup\pmb{a}$ and $E=\{\pmb{b}\}$.
  Then $A^*E=\{\pmb{a}\}^*\{\pmb{b}\}=\{\pmb{a}^n\pmb{b}\mid n\ge 0\}$.
  But $X=\Sigma^*$ also satisfies $X=A\cdot X\cup E$: indeed,
  \[
    A\cdot\Sigma^*\cup\{\pmb{b}\}
    =(\{\lambda,\pmb{a}\}\cdot\Sigma^*)\cup\{\pmb{b}\}
    =\Sigma^*\cup\{\pmb{b}\}
    =\Sigma^*.
  \]
  So $\Sigma^*$ is also a solution, but $\Sigma^*\neq A^*E$.  Hence the solution is not unique.

  \item \textbf{$A^*E$ as unique solution even if $\lambda\in A$.}

  Let the alphabet be $\{\pmb{a}\}$, let $A=\{\lambda,\pmb{a}\}$, and let
  $E=\pmb{a}^*$.  Then
  \[
    A^*E=(\epsilon\cup\pmb{a})^*\pmb{a}^*=\pmb{a}^*.
  \]
  The equation becomes
  \[
    X=E\cup A\cdot X=\pmb{a}^*\cup(\epsilon\cup\pmb{a})X.
  \]
  Any solution must contain $E=\pmb{a}^*$, but $\pmb{a}^*$ is already the whole
  universe over this alphabet.  Therefore the only possible solution is
  $X=\pmb{a}^*=A^*E$.
\end{enumerate}

%% 6.18
\item \textbf{Solve the system of Exercise~6.18.}

The equations are
\[
X_0=(\pmb{0}\cup\pmb{1})X_1,
\qquad
X_1=\epsilon\cup\pmb{1}X_0\cup\pmb{0}X_1.
\]
Substitute the first equation into the second:
\[
X_1=\epsilon\cup\pmb{1}(\pmb{0}\cup\pmb{1})X_1\cup\pmb{0}X_1
=\epsilon\cup(\pmb{10}\cup\pmb{11}\cup\pmb{0})X_1.
\]
Since $\lambda\not in(\pmb{10}\cup\pmb{11}\cup\pmb{0})$, Theorem~6.1 gives
\[
X_1=(\pmb{10}\cup\pmb{11}\cup\pmb{0})^*.
\]
Therefore
\[
X_0=(\pmb{0}\cup\pmb{1})X_1
=(\pmb{0}\cup\pmb{1})(\pmb{10}\cup\pmb{11}\cup\pmb{0})^*.
\]
These are the terminal sets for the corresponding two-state DFA.
In fact, this is the same state-equation system one obtains from the DFA in
Example~3.4, where $X_0$ and $X_1$ are the two terminal sets.

%% 6.19
\item \textbf{Solve for $X_1,X_2,X_3$ from the equations.}

Simplifying (noting $\emptyset\cdot X_i=\emptyset$):
\begin{align*}
  X_1 &= (\pmb{0}\cup\pmb{1})X_2 \\
  X_2 &= \epsilon\cup\pmb{0}X_1\cup\pmb{1}X_2 \\
  X_3 &= (\pmb{0}\cup\pmb{1})X_2
\end{align*}
Note $X_3=X_1$.

\begin{enumerate}[(a)]
  \item \textbf{Eliminate $X_3$, then $X_2$:}

  Since $X_3=X_1$, eliminate $X_3$.  From $X_2=\epsilon\cup\pmb{0}X_1\cup\pmb{1}X_2$: by Theorem~6.1 ($\lambda\not in\{\pmb{1}\}$): $X_2=\pmb{1}^*(\epsilon\cup\pmb{0}X_1)=\pmb{1}^*\cup\pmb{1}^*\pmb{0}X_1$.
  Substitute into $X_1=(\pmb{0}\cup\pmb{1})X_2=(\pmb{0}\cup\pmb{1})(\pmb{1}^*\cup\pmb{1}^*\pmb{0}X_1)=(\pmb{0}\cup\pmb{1})\pmb{1}^*\cup(\pmb{0}\cup\pmb{1})\pmb{1}^*\pmb{0}X_1$.
  By Theorem~6.1: $X_1=((\pmb{0}\cup\pmb{1})\pmb{1}^*\pmb{0})^*(\pmb{0}\cup\pmb{1})\pmb{1}^*$.
  Then
  \[
    X_2=\pmb{1}^*\cup\pmb{1}^*\pmb{0}((\pmb{0}\cup\pmb{1})\pmb{1}^*\pmb{0})^*(\pmb{0}\cup\pmb{1})\pmb{1}^*,
    \qquad
    X_3=X_1.
  \]

  \item \textbf{Eliminate $X_3$, then $X_1$:}

  Same $X_3=X_1$.  From $X_1=(\pmb{0}\cup\pmb{1})X_2$, substitute into $X_2$:
  $X_2=\epsilon\cup\pmb{0}(\pmb{0}\cup\pmb{1})X_2\cup\pmb{1}X_2=\epsilon\cup(\pmb{0}(\pmb{0}\cup\pmb{1})\cup\pmb{1})X_2$.
  Let $A=\pmb{0}(\pmb{0}\cup\pmb{1})\cup\pmb{1}=\pmb{00}\cup\pmb{01}\cup\pmb{1}$; $\lambda\not in A$.
  $X_2=A^*$.
  Then $X_1=(\pmb{0}\cup\pmb{1})A^*$ and $X_3=X_1$.

  \item \textbf{Comparison:} Both solutions are equivalent: $(\pmb{0}\cup\pmb{1})\pmb{1}^*(\pmb{0}(\pmb{0}\cup\pmb{1})\pmb{1}^*)^*=(\pmb{0}\cup\pmb{1})(\pmb{00}\cup\pmb{01}\cup\pmb{1})^*$.  The second form is more concise.
\end{enumerate}

%% 6.20
\item \textbf{Prove Lemma~6.2: every regular expression describes a FAD language.}

\begin{proof}
Let $P(m)$: ``Every regular expression with $\le m$ operators represents a FAD language.''

\emph{Base ($m=0$):}  No operators: the expression is $\emptyset$, $\epsilon$, or a single letter $\pmb{a}\in\Sigma$.  These represent $\emptyset$, $\{\lambda\}$, and $\{\pmb{a}\}$ respectively---all finite, hence FAD.

\emph{Inductive step:}  Assume $P(m)$.  Any regular expression with $m+1$ operators is either $R_1\cup R_2$, $R_1\cdot R_2$, or $R_1^*$, where $R_1$ and $R_2$ each have $\le m$ operators.  By $P(m)$, $L(R_1),L(R_2)\in\DSIG$.  Then:
\begin{itemize}
  \item $L(R_1\cup R_2)=L(R_1)\cup L(R_2)\in\DSIG$ (Theorem~5.2).
  \item $L(R_1\cdot R_2)=L(R_1)\cdot L(R_2)\in\DSIG$ (Theorem~5.4).
  \item $L(R_1^*)=L(R_1)^*\in\DSIG$ (Theorem~5.5).
\end{itemize}
So $P(m+1)$ holds.  By induction, every regular expression represents a FAD language. \qed
\end{proof}

%% 6.21
\item \textbf{All solutions to $X=X\cup\{\pmb{b}\}$ over $\{\pmb{a},\pmb{b},\pmb{c}\}^*$.}

$X=X\cup\{\pmb{b}\}$ means exactly that $\{\pmb{b}\}\subseteq X$.  Hence the
solutions are precisely the languages $X$ with $\pmb{b}\in X$.  Every such
language is a solution, and the smallest one is $\{\pmb{b}\}$.

%% 6.22
\item \textbf{Prove $A^*E=E\cup A\cdot(A^*E)$ for any languages $A$ and $E$.}

\begin{proof}
$A^*E=(\epsilon\cup A\cup A^2\cup\cdots)E=\epsilon E\cup AE\cup A^2E\cup\cdots=E\cup A(E\cup AE\cup A^2E\cup\cdots)=E\cup A(A^*E)$. \qed
\end{proof}

%% 6.23
\item \textbf{Regular expression for $(\pmb{ab}\cup\pmb{b})^*\pmb{a}\cap(\pmb{ba}\cup\pmb{a})^*$.}

Both are FAD; their intersection is FAD (Lemma~5.1).  We use the product DFA approach.

$L_1=(\pmb{ab}\cup\pmb{b})^*\pmb{a}$: strings over $\{\pmb{a},\pmb{b}\}$ of the form $w\pmb{a}$ where $w\in(\pmb{ab}\cup\pmb{b})^*$.
$L_2=(\pmb{ba}\cup\pmb{a})^*$: strings over $\{\pmb{a},\pmb{b}\}$ that are concatenations of $\pmb{ba}$ or $\pmb{a}$.

$L_1$ ends in $\pmb{a}$.  $L_2=(\pmb{ba}\cup\pmb{a})^*$ consists of all strings that decompose as concatenations of $\pmb{a}$ and $\pmb{ba}$.

$L_2=(\pmb{ba}\cup\pmb{a})^*$ means every $\pmb{b}$ is immediately followed by $\pmb{a}$.  A word in $L_1=(\pmb{ab}\cup\pmb{b})^*\pmb{a}$ has prefix in $(\pmb{ab}\cup\pmb{b})^*$, so every $\pmb{a}$ in the prefix is immediately followed by $\pmb{b}$.  Combining both constraints on the prefix atoms: each atom (whether $\pmb{ab}$ or $\pmb{b}$) must be immediately followed by something starting with $\pmb{a}$ (since the atom ends in $\pmb{b}$, and every $\pmb{b}$ must be followed by $\pmb{a}$).  Consequently, $\pmb{b}$-atoms can appear only once, at the start (any $\pmb{b}$-atom followed by another $\pmb{b}$-atom would put $\pmb{b}$ before $\pmb{b}$, violating $L_2$); and $\pmb{ab}$-atoms must be followed by $\pmb{ab}$-atoms or the final $\pmb{a}$.  The intersection is therefore:
\[
  L_1\cap L_2 = (\epsilon\cup\pmb{b})(\pmb{ab})^*\pmb{a}
\]
Verification: $(\epsilon\cup\pmb{b})(\pmb{ab})^k\pmb{a}$ gives words $\pmb{a}, \pmb{aba}, \pmb{ababa},\ldots$ and $\pmb{ba}, \pmb{baba}, \pmb{bababa},\ldots$, all of which belong to both $L_1$ and $L_2$; and any word in $L_1\cap L_2$ has this form by the above analysis.

%% 6.24
\item \textbf{Algorithm for intersection of two regular expressions.}

Given regular expressions $R_1$ and $R_2$ (each describing a FAD language):
\begin{enumerate}[1.]
  \item Convert $R_1$ to an NDFA $A_1$ using the constructions of Chapter~5.
  \item Convert $R_2$ to an NDFA $A_2$ similarly.
  \item Apply the subset construction to get DFAs $A_1^d$ and $A_2^d$.
  \item Build the product DFA $A^\cap$ (Lemma~5.1) with $F^\cap=F_1\times F_2$.
  \item Use the state-equation method of Theorems~6.3 and~6.2 to convert $A^\cap$ to a regular expression for $L(A_1)\cap L(A_2)$.
\end{enumerate}

%% 6.25
\item \textbf{Verify that $X_1=(\pmb{a}\cup\pmb{bb})^*$ and $X_2=\pmb{b}\cdot(\pmb{a}\cup\pmb{bb})^*$ satisfy:}
\begin{align*}
  X_1 &= \epsilon\cup\pmb{a}\cdot X_1\cup\pmb{b}\cdot X_2 \\
  X_2 &= \pmb{b}\cdot X_1
\end{align*}

\textbf{Verification of $X_2=\pmb{b}\cdot X_1$:}
\[
  \pmb{b}\cdot X_1 = \pmb{b}\cdot(\pmb{a}\cup\pmb{bb})^* = X_2. \quad\checkmark
\]

\textbf{Verification of $X_1=\epsilon\cup\pmb{a}\cdot X_1\cup\pmb{b}\cdot X_2$:}
\begin{align*}
  \epsilon\cup\pmb{a}\cdot X_1\cup\pmb{b}\cdot X_2
    &= \epsilon\cup\pmb{a}(\pmb{a}\cup\pmb{bb})^*\cup\pmb{b}\cdot\pmb{b}(\pmb{a}\cup\pmb{bb})^* \\
    &= \epsilon\cup(\pmb{a}\cup\pmb{bb})\cdot(\pmb{a}\cup\pmb{bb})^* \\
    &= (\pmb{a}\cup\pmb{bb})^* = X_1. \quad\checkmark
\end{align*}

%% 6.26
\item \begin{enumerate}[(a)]
  \item \textbf{$L(D)$ for Example~6.10:}  Example~6.10 gives
  $L(B)=X_1=(\pmb{a}\cup\pmb{bb})^*$ and
  $L(C)=X_2=\pmb{b}(\pmb{a}\cup\pmb{bb})^*$.  Since $D$ has both corresponding
  start states, its language is
  \[
    L(D)=X_1\cup X_2
    =(\epsilon\cup\pmb{b})(\pmb{a}\cup\pmb{bb})^*.
  \]

  \item \textbf{General formula:}  For a machine $A$ with start states $s_{i_1},\ldots,s_{i_m}$:
  \[
    L(A) = X_{i_1} \cup X_{i_2} \cup \cdots \cup X_{i_m}
  \]
  where each $X_j$ is the solution to the state equation for state $s_j$.
\end{enumerate}

%% 6.27
\item \textbf{Regular expression for the complement of $(\pmb{ab}\cup\pmb{b})^*\pmb{a}$.}

The language $(\pmb{ab}\cup\pmb{b})^*\pmb{a}$ is exactly the set of words over
$\{\pmb{a},\pmb{b}\}$ that end in $\pmb{a}$ and contain no $\pmb{aa}$ as a
substring.  Therefore its complement is the set of words that either end in
$\pmb{b}$ (or are $\lambda$) or contain $\pmb{aa}$.  One regular expression is
\[
  \epsilon \cup (\pmb{a}\cup\pmb{b})^*\pmb{b}
  \cup (\pmb{a}\cup\pmb{b})^*\pmb{aa}(\pmb{a}\cup\pmb{b})^*.
\]

%% 6.28
\item \textbf{Algorithm to complement a regular expression.}

Given regular expression $R$:
\begin{enumerate}[1.]
  \item Convert $R$ to a DFA $A$ (via NDFA $\to$ subset construction).
  \item Complement $A$ by swapping $F$ and $S\setminus F$.
  \item Use the state-equation method of Theorems~6.3 and~6.2 to convert the complemented DFA back to a regular expression.
\end{enumerate}

%% 6.29
\item \textbf{$\RSIG$ is closed under $E(L)$ using regular expressions.}

Here
\[
E(L)=\{z\mid(\exists y\in\Sigma^+)(\exists x\in L)\ z=yx\}=\Sigma^+\cdot L.
\]
If $L=L(R)$ and $\Sigma=\{\pmb{a},\pmb{b},\pmb{c}\}$, then
\[
E(L)=L\Bigl((\pmb{a}\cup\pmb{b}\cup\pmb{c})(\pmb{a}\cup\pmb{b}\cup\pmb{c})^*R\Bigr),
\]
which is again a regular expression.  Hence \RSIG\ is closed under $E$.

%% 6.30
\item \textbf{$\RSIG$ is closed under $B(L)$.}

$B(L)=L\cdot\Sigma^*$ (words starting with a word from $L$).  If $L=L(R)$, then $B(L)=L(R\cdot\Sigma^*)=L(R\cdot(\pmb{a}\cup\pmb{b}\cup\pmb{c})^*)$.  This is a regular expression, so $B(L)\in\RSIG$.

%% 6.31
\item \textbf{$\RSIG$ is closed under $M(L)$.}

$M(L)=L\cdot\Sigma^+$.  If $L=L(R)$, then $M(L)=L(R\cdot(\pmb{a}\cup\pmb{b}\cup\pmb{c})^+)=L(R\cdot(\pmb{a}\cup\pmb{b}\cup\pmb{c})\cdot(\pmb{a}\cup\pmb{b}\cup\pmb{c})^*)$.  Regular expression, hence $M(L)\in\RSIG$.

%% 6.32
\item \begin{enumerate}[(a)]
  \item \textbf{No unique solution to the given system:}

  The equations:
  $X_1=\pmb{b}\cup\epsilon\cdot X_1\cup\pmb{a}\cdot X_2=\pmb{b}\cup X_1\cup\pmb{a}X_2$.
  Since $X_1\subseteq X_1\cup\pmb{b}\cup\pmb{a}X_2$, the equation $X_1=\pmb{b}\cup X_1\cup\pmb{a}X_2$ says $X_1\supseteq\pmb{b}\cup\pmb{a}X_2$.  Any $X_1$ that contains $\pmb{b}$ and $\pmb{a}X_2$ is a solution; there are infinitely many such $X_1$.

  $X_2=\pmb{c}\cup\emptyset\cdot X_1\cup\epsilon\cdot X_2=\pmb{c}\cup X_2$.  Similarly, any $X_2\supseteq\{\pmb{c}\}$ is a solution.

  \item \textbf{No contradiction with Theorem~6.2:}  Theorem~6.2 assumes
  $\lambda\not in A_{ij}$ for \emph{every} coefficient in the system.  Here
  $A_{11}=\epsilon$ and $A_{22}=\epsilon$, so the hypothesis fails.  Therefore
  Theorem~6.2 does not apply, and there is no contradiction.
\end{enumerate}

%% 6.33
\item \textbf{Solve for $X_0$ and $X_1$.}

\begin{align*}
  X_0 &= \pmb{0}^*\pmb{1} \cup (\pmb{10})^*X_0 \cup \pmb{0}(\pmb{0}\cup\pmb{1})X_1 \\
  X_1 &= \epsilon \cup \pmb{1}^*\pmb{01}X_0 \cup \pmb{0}X_1
\end{align*}

From the $X_1$ equation (by Theorem~6.1, $\lambda\not in\{\pmb{0}\}$):
$X_1=\pmb{0}^*(\epsilon\cup\pmb{1}^*\pmb{01}X_0)=\pmb{0}^*\cup\pmb{0}^*\pmb{1}^*\pmb{01}X_0$.

Substitute into $X_0$:
$X_0=\pmb{0}^*\pmb{1}\cup(\pmb{10})^*X_0\cup\pmb{0}(\pmb{0}\cup\pmb{1})(\pmb{0}^*\cup\pmb{0}^*\pmb{1}^*\pmb{01}X_0)$.

Let $C=\pmb{0}(\pmb{0}\cup\pmb{1})\pmb{0}^*\pmb{1}^*\pmb{01}$. Then:
$X_0=((\pmb{10})^*\cup C)X_0\cup(\pmb{0}^*\pmb{1}\cup\pmb{0}(\pmb{0}\cup\pmb{1})\pmb{0}^*)$.

Since $(\pmb{10})^*$ contains $\lambda$, Theorem~6.1 does not guarantee a
unique solution here.  Following the chapter's discussion of ambiguous
equations, the least solution is
\[
  X_0=((\pmb{10})^*\cup C)^*(\pmb{0}^*\pmb{1}\cup\pmb{0}(\pmb{0}\cup\pmb{1})\pmb{0}^*).
\]
Back-substitution then gives
\[
  X_1=
  \pmb{0}^*\cup
  \pmb{0}^*\pmb{1}^*\pmb{01}
  ((\pmb{10})^*\cup C)^*(\pmb{0}^*\pmb{1}\cup\pmb{0}(\pmb{0}\cup\pmb{1})\pmb{0}^*).
\]

%% 6.34
\item \begin{enumerate}[(a)]
  \item Words ending with $\pmb{c}$ with no $\pmb{aa}$, $\pmb{bb}$, $\pmb{cc}$ substrings:

  Build the DFA with states $q_0$ (start), $q_a$ (last letter $\pmb{a}$), $q_b$ (last letter $\pmb{b}$), $q_c$ (last letter $\pmb{c}$, accepting), and dead state $q_d$; transitions $q_x\xrightarrow{x}q_d$ and $q_x\xrightarrow{y}q_y$ for $y\neq x$.  Write state equations for strings reaching $q_c$:
  \begin{align*}
    X_a &= \pmb{b} X_b \cup \pmb{c} X_c \\
    X_b &= \pmb{a} X_a \cup \pmb{c} X_c \\
    X_c &= \epsilon \cup \pmb{a} X_a \cup \pmb{b} X_b
  \end{align*}
  Eliminate $X_a$: substituting $X_b = \pmb{a}X_a\cup\pmb{c}X_c$ into $X_a$ gives $X_a = \pmb{ba}X_a\cup(\pmb{bc}\cup\pmb{c})X_c$, so $X_a=(\pmb{ba})^*(\pmb{b}\cup\epsilon)\pmb{c} X_c$.  Then $X_b = \pmb{a}(\pmb{ba})^*(\pmb{b}\cup\epsilon)\pmb{c} X_c\cup\pmb{c} X_c$.  Substituting into $X_c$:
  \[
    X_c = \epsilon \cup \bigl((\pmb{a}\cup\pmb{ba})(\pmb{ba})^*(\pmb{b}\cup\epsilon)\pmb{c}\cup\pmb{bc}\bigr)X_c
  \]
  Since $\lambda$ is absent from the coefficient, $X_c = A^*$ where
  \[
    A = (\pmb{a}\cup\pmb{ba})(\pmb{ba})^*(\pmb{b}\cup\epsilon)\pmb{c}\cup\pmb{bc}.
  \]
  The language is $L = \pmb{a} X_a\cup\pmb{b} X_b\cup\pmb{c} X_c = (\pmb{a}(\pmb{ba})^*(\pmb{b}\cup\epsilon)\pmb{c} \cup \pmb{b}(\pmb{a}(\pmb{ba})^*(\pmb{b}\cup\epsilon)\pmb{c}\cup\pmb{c})\cup\pmb{c}) A^*$.

  \item Words beginning with $\pmb{c}$ with no $\pmb{aa}$, $\pmb{bb}$, $\pmb{cc}$ substrings:
  Let $X_c$ be the set of suffixes that may follow an initial $\pmb{c}$ while
  preserving the ``no equal consecutive letters'' condition.  Solving the
  corresponding state equations gives
  \[
    X_c=
    \bigl(((\pmb{a}\cup\pmb{ba})(\pmb{ba})^*(\epsilon\cup\pmb{b})\cup\pmb{b})\pmb{c}\bigr)^*
    (\epsilon\cup\pmb{a})(\pmb{ba})^*(\epsilon\cup\pmb{b}).
  \]
  Therefore one regular expression for the language is
  \[
    \pmb{c}\bigl(((\pmb{a}\cup\pmb{ba})(\pmb{ba})^*(\epsilon\cup\pmb{b})\cup\pmb{b})\pmb{c}\bigr)^*
    (\epsilon\cup\pmb{a})(\pmb{ba})^*(\epsilon\cup\pmb{b}).
  \]
\end{enumerate}

%% 6.35
\item \begin{enumerate}[(a)]
  \item Words with at most two $\pmb{c}$s:
  \[
    (\pmb{a}\cup\pmb{b})^*\cup(\pmb{a}\cup\pmb{b})^*\pmb{c}(\pmb{a}\cup\pmb{b})^*\cup(\pmb{a}\cup\pmb{b})^*\pmb{c}(\pmb{a}\cup\pmb{b})^*\pmb{c}(\pmb{a}\cup\pmb{b})^*
  \]

  \item Words that do \emph{not} have exactly one $\pmb{c}$ (i.e., zero or $\ge 2$ $\pmb{c}$s):
  \[
    (\pmb{a}\cup\pmb{b})^*
    \cup
    (\pmb{a}\cup\pmb{b})^*\pmb{c}(\pmb{a}\cup\pmb{b})^*\pmb{c}(\pmb{a}\cup\pmb{b}\cup\pmb{c})^*
  \]
\end{enumerate}

%% 6.36
\item \textbf{Reversal of regular expressions.}

\begin{enumerate}[(a)]
  \item Equivalences:
  \begin{itemize}
    \item $(R_1\cup R_2)^r=R_1^r\cup R_2^r$
    \item $(R_1\cdot R_2)^r=R_2^r\cdot R_1^r$
    \item $(R^*)^r=(R^r)^*$
    \item $\emptyset^r=\emptyset$
    \item $\epsilon^r=\epsilon$
    \item $\pmb{a}^r=\pmb{a}$ for each $\pmb{a}\in\Sigma$
  \end{itemize}

  \item To generate a regular expression for $L^r$ from $R$: apply the reversal operator recursively to $R$ using the rules above.

  \item \textbf{Proof by induction on the number of operators $m$ in $R$:}

  \emph{Base ($m=0$):}  $R\in\{\emptyset,\epsilon,\pmb{a}\}$.  $L(R)^r=L(R^r)$ by direct check.

  \emph{Step:}  Assume the result for all expressions with $<m$ operators.  If $R=R_1\cup R_2$: $(L(R_1)\cup L(R_2))^r=L(R_1)^r\cup L(R_2)^r=L(R_1^r)\cup L(R_2^r)=L(R_1^r\cup R_2^r)=L(R^r)$.  Similarly for $\cdot$ and $*$.  \qed

  \item Since every FAD language has a regular expression (Corollary~6.2), and
  the reversal of any regular expression (via the rules above) is also a
  regular expression, $\RSIG$ is closed under $r$.
\end{enumerate}

%% 6.37
\item \textbf{Complete details of Theorem~6.4.}

Theorem~6.4 is the left-linear analog of Theorem~6.2.

\begin{proof}
Focus on the $m$th equation:
\[
  Y_m=I_m\cup Y_1B_{m1}\cup\cdots\cup Y_{m-1}B_{m(m-1)}\cup Y_mB_{mm}
      \cup Y_{m+1}B_{m(m+1)}\cup\cdots\cup Y_nB_{mn}.
\]
Treating every term except $Y_mB_{mm}$ as a constant, the mirror-image of the
argument in Theorem~6.1 gives
\[
  Y_m=
  (I_m\cup Y_1B_{m1}\cup\cdots\cup Y_{m-1}B_{m(m-1)}
  \cup Y_{m+1}B_{m(m+1)}\cup\cdots\cup Y_nB_{mn})B_{mm}^*.
\]
This is part (b).

Substitute this expression for $Y_m$ into each remaining equation
\[
  Y_k=I_k\cup Y_1B_{k1}\cup\cdots\cup Y_mB_{km}\cup\cdots\cup Y_nB_{kn}
  \qquad (k\ne m),
\]
and distribute $B_{km}$ across the substituted expression.  Collecting the
constant term and the coefficients of each $Y_j$ yields
\[
  \hat{I}_k=I_k\cup(I_mB_{mm}^*B_{km})
\]
and
\[
  \hat{B}_{kj}=B_{kj}\cup(B_{mj}B_{mm}^*B_{km})
  \qquad (j\ne m),
\]
which is exactly the reduced system in part (a).

For a single left-linear equation
\[
  Y_1=I_1\cup Y_1B_{11},
\]
the same mirror-image argument gives the unique solution
\[
  Y_1=I_1B_{11}^*.
\]
This is part (c).  Hence Theorem~6.4 follows. \qed
\end{proof}

%% 6.38
\item \begin{enumerate}[(a)]
  \item Words with no $\pmb{b}$ immediately preceded by $\pmb{a}$ over $\{\pmb{a},\pmb{b},\pmb{c}\}$:

  This means exactly that $\pmb{ab}$ never occurs.  Every maximal block of
  $\pmb{a}$s must therefore be followed by $\pmb{c}$ or occur at the end of the
  word.  Hence one regular expression is
  \[
    (\pmb{b}\cup\pmb{c}\cup\pmb{a}\pmb{a}^*\pmb{c})^*(\epsilon\cup\pmb{a}\pmb{a}^*).
  \]

  \item Words with exactly two $\pmb{c}$s and no $\pmb{ab}$:

  Over $\{\pmb{a},\pmb{b}\}$, avoiding $\pmb{ab}$ means that once an
  $\pmb{a}$ appears, only $\pmb{a}$s may follow; hence the $\pmb{a}/\pmb{b}$
  portions are exactly $\pmb{b}^*\pmb{a}^*$.  Therefore the desired language is
  \[
    \pmb{b}^*\pmb{a}^*\pmb{c}\pmb{b}^*\pmb{a}^*\pmb{c}\pmb{b}^*\pmb{a}^*.
  \]
\end{enumerate}

%% 6.39
\item \begin{enumerate}[(a)]
  \item Words with no $\pmb{b}$ immediately preceded by $\pmb{c}$ (no $\pmb{cb}$ substring):

  This means exactly that every maximal block of $\pmb{c}$s is followed by
  $\pmb{a}$ or occurs at the end of the word.  Hence one regular expression is
  \[
    (\pmb{a}\cup\pmb{b}\cup\pmb{c}\pmb{c}^*\pmb{a})^*(\epsilon\cup\pmb{c}\pmb{c}^*)
  \]

  \item Words with exactly one $\pmb{c}$ and no $\pmb{cb}$:

  Exactly one $\pmb{c}$, and that $\pmb{c}$ is not immediately followed by $\pmb{b}$:
  \[
    (\pmb{a}\cup\pmb{b})^*\pmb{c}(\pmb{a}\cup\epsilon)(\pmb{a}\cup\pmb{b})^*
    \quad\text{where the $\pmb{c}$ is followed by $\pmb{a}$ or end}
  \]
  More precisely:
  \[
    (\pmb{a}\cup\pmb{b})^*\cdot\pmb{c}\cdot((\pmb{a}(\pmb{a}\cup\pmb{b})^*)\cup\epsilon)
  \]
\end{enumerate}

%% 6.40
\item \begin{enumerate}[(a)]
  \item For Figure~6.11 the right-linear equations are
  \begin{align*}
    X_1 &= \emptyset \cup \pmb{a}X_1 \cup \pmb{a}X_2,\\
    X_2 &= \epsilon \cup (\pmb{a}\cup\pmb{b})X_1.
  \end{align*}

  \item Substitute the second equation into the first:
  \[
  X_1=\pmb{a}X_1\cup\pmb{a}(\epsilon\cup(\pmb{a}\cup\pmb{b})X_1)
     =\pmb{a}\cup(\pmb{a}\cup\pmb{aa}\cup\pmb{ab})X_1.
  \]
  Hence
  \[
  X_1=(\pmb{a}\cup\pmb{aa}\cup\pmb{ab})^*\pmb{a},
  \qquad
  X_2=\epsilon\cup(\pmb{a}\cup\pmb{b})(\pmb{a}\cup\pmb{aa}\cup\pmb{ab})^*\pmb{a}.
  \]

  \item Therefore the NDFA accepts
  \[
    L=(\pmb{a}\cup\pmb{aa}\cup\pmb{ab})^*\pmb{a}.
  \]

  \item The left-linear equations are
  \begin{align*}
    Y_1 &= \epsilon \cup Y_1\pmb{a} \cup Y_2(\pmb{a}\cup\pmb{b}),\\
    Y_2 &= Y_1\pmb{a}.
  \end{align*}
  Substituting gives
  \[
  Y_1=\epsilon\cup Y_1(\pmb{a}\cup\pmb{aa}\cup\pmb{ab}),
  \]
  so
  \[
  Y_1=(\pmb{a}\cup\pmb{aa}\cup\pmb{ab})^*,
  \qquad
  Y_2=(\pmb{a}\cup\pmb{aa}\cup\pmb{ab})^*\pmb{a}.
  \]
\end{enumerate}

%% 6.41
\item \begin{enumerate}[(a)]
  \item For Figure~6.12 the right-linear equations are
  \begin{align*}
    X_0 &= \pmb{a}X_1,\\
    X_1 &= \pmb{a}X_1 \cup \pmb{a}X_2,\\
    X_2 &= \epsilon \cup \pmb{b}X_2 \cup \pmb{a}X_3,\\
    X_3 &= \pmb{b}X_0 \cup (\pmb{a}\cup\pmb{b})X_3 .
  \end{align*}

  \item Solving from the end:
  \[
  X_3=(\pmb{a}\cup\pmb{b})^*\pmb{b}X_0,
  \]
  \[
  X_2=\pmb{b}^*(\epsilon\cup\pmb{a}X_3)
      =\pmb{b}^*\cup\pmb{b}^*\pmb{a}(\pmb{a}\cup\pmb{b})^*\pmb{b}X_0,
  \]
  \[
  X_1=\pmb{a}^*\pmb{a}X_2,
  \qquad
  X_0=\pmb{a}\pmb{a}^*\pmb{a}X_2.
  \]
  Therefore
  \[
  X_0=\pmb{a}\pmb{a}^*\pmb{a}\pmb{b}^*
     \cup \pmb{a}\pmb{a}^*\pmb{a}\pmb{b}^*\pmb{a}(\pmb{a}\cup\pmb{b})^*\pmb{b}X_0,
  \]
  so
  \[
  X_0=
  \bigl(\pmb{a}\pmb{a}^*\pmb{a}\pmb{b}^*\pmb{a}(\pmb{a}\cup\pmb{b})^*\pmb{b}\bigr)^*
  \pmb{a}\pmb{a}^*\pmb{a}\pmb{b}^*.
  \]

  \item Hence the language of the automaton is
  \[
  L=
  \bigl(\pmb{a}\pmb{a}^*\pmb{a}\pmb{b}^*\pmb{a}(\pmb{a}\cup\pmb{b})^*\pmb{b}\bigr)^*
  \pmb{a}\pmb{a}^*\pmb{a}\pmb{b}^*.
  \]

  \item The left-linear equations are
  \begin{align*}
    Y_0 &= \epsilon \cup Y_3\pmb{b},\\
    Y_1 &= Y_0\pmb{a} \cup Y_1\pmb{a},\\
    Y_2 &= Y_1\pmb{a} \cup Y_2\pmb{b},\\
    Y_3 &= Y_2\pmb{a} \cup Y_3(\pmb{a}\cup\pmb{b}).
  \end{align*}
  Solving gives
  \[
  Y_0=
  \bigl(\pmb{a}\pmb{a}^*\pmb{a}\pmb{b}^*\pmb{a}(\pmb{a}\cup\pmb{b})^*\pmb{b}\bigr)^*,
  \]
  \[
  Y_1=Y_0\pmb{a}\pmb{a}^*,
  \qquad
  Y_2=Y_0\pmb{a}\pmb{a}^*\pmb{a}\pmb{b}^*,
  \qquad
  Y_3=Y_0\pmb{a}\pmb{a}^*\pmb{a}\pmb{b}^*\pmb{a}(\pmb{a}\cup\pmb{b})^*.
  \]
\end{enumerate}

%% 6.42
\item \begin{enumerate}[(a)]
  \item For Figure~6.13 the right-linear equations are
  \begin{align*}
    X_0 &= \pmb{0}X_0 \cup \pmb{1}X_0 \cup \pmb{0}X_1,\\
    X_1 &= \pmb{1}X_2,\\
    X_2 &= \pmb{0}X_3,\\
    X_3 &= \pmb{0}X_4,\\
    X_4 &= \pmb{1}X_5,\\
    X_5 &= \pmb{0}X_6,\\
    X_6 &= \epsilon \cup \pmb{0}X_6 \cup \pmb{1}X_6 .
  \end{align*}

  \item These equations collapse immediately:
  \[
  X_6=(\pmb{0}\cup\pmb{1})^*,
  \quad
  X_5=\pmb{0}(\pmb{0}\cup\pmb{1})^*,
  \quad
  X_4=\pmb{10}(\pmb{0}\cup\pmb{1})^*,
  \]
  \[
  X_3=\pmb{010}(\pmb{0}\cup\pmb{1})^*,
  \quad
  X_2=\pmb{0010}(\pmb{0}\cup\pmb{1})^*,
  \quad
  X_1=\pmb{10010}(\pmb{0}\cup\pmb{1})^*.
  \]
  Hence
  \[
  X_0=(\pmb{0}\cup\pmb{1})X_0 \cup \pmb{010010}(\pmb{0}\cup\pmb{1})^*,
  \]
  so
  \[
  X_0=(\pmb{0}\cup\pmb{1})^*\pmb{010010}(\pmb{0}\cup\pmb{1})^*.
  \]

  \item Therefore the language is all binary strings containing
  $\pmb{010010}$ as a substring:
  \[
  L=(\pmb{0}\cup\pmb{1})^*\pmb{010010}(\pmb{0}\cup\pmb{1})^*.
  \]

  \item The left-linear equations are
  \begin{align*}
    Y_0 &= \epsilon \cup Y_0\pmb{0} \cup Y_0\pmb{1},\\
    Y_1 &= Y_0\pmb{0},\\
    Y_2 &= Y_1\pmb{1},\\
    Y_3 &= Y_2\pmb{0},\\
    Y_4 &= Y_3\pmb{0},\\
    Y_5 &= Y_4\pmb{1},\\
    Y_6 &= Y_5\pmb{0} \cup Y_6\pmb{0} \cup Y_6\pmb{1}.
  \end{align*}
  Solving gives
  \[
  Y_0=(\pmb{0}\cup\pmb{1})^*,
  \quad
  Y_1=(\pmb{0}\cup\pmb{1})^*\pmb{0},
  \quad
  Y_2=(\pmb{0}\cup\pmb{1})^*\pmb{01},
  \]
  \[
  Y_3=(\pmb{0}\cup\pmb{1})^*\pmb{010},
  \quad
  Y_4=(\pmb{0}\cup\pmb{1})^*\pmb{0100},
  \quad
  Y_5=(\pmb{0}\cup\pmb{1})^*\pmb{01001},
  \]
  and
  \[
  Y_6=(\pmb{0}\cup\pmb{1})^*\pmb{010010}(\pmb{0}\cup\pmb{1})^*.
  \]
\end{enumerate}

%% 6.43
\item \textbf{Prove: for any $A$, $E$, $Y$, if $E\subseteq Y$ then $A\cdot E\subseteq A\cdot Y$.}

\begin{proof}
Let $w\in A\cdot E$.  Then $w=xy$ with $x\in A$ and $y\in E$.  Since $E\subseteq Y$, $y\in Y$.  Hence $w=xy\in A\cdot Y$. \qed
\end{proof}

%% 6.44
\item \begin{enumerate}[(a)]
  \item \textbf{Image of \RSIG\ under substitution $\bar{s}$ is in $\mathfrak{R}_\Gamma$.}

  \begin{proof}
  Interpreting $s$ as a regular set substitution in the sense of
  Definition~6.5, this is exactly the induction omitted in Theorem~6.6.
  For a regular expression $R$ over $\Sigma$, define $\bar{s}(R)$ by:
  $\bar{s}(\pmb{a})=s(\pmb{a})\in\mathfrak{R}_\Gamma$ (by assumption),
  $\bar{s}(\emptyset)=\emptyset$, $\bar{s}(\epsilon)=\{\lambda\}$,
  $\bar{s}(R_1\cup R_2)=\bar{s}(R_1)\cup\bar{s}(R_2)\in\mathfrak{R}_\Gamma$ (closed under union),
  $\bar{s}(R_1\cdot R_2)=\bar{s}(R_1)\cdot\bar{s}(R_2)\in\mathfrak{R}_\Gamma$ (closed under concatenation),
  $\bar{s}(R^*)=\bar{s}(R)^*\in\mathfrak{R}_\Gamma$ (closed under Kleene closure).
  By induction on the structure of $R$, $\bar{s}(R)\in\mathfrak{R}_\Gamma$. \qed
  \end{proof}

  \item \textbf{Image of \NSIG\ under $\bar{s}$ need not be in $\mathfrak{N}_\Gamma$.}

  Let
  \[
  L=\{x\in\{\pmb{a},\pmb{b}\}^*\mid |x|_{\pmb{a}}=|x|_{\pmb{b}}\}\in\NSIG.
  \]
  Define the substitution by
  $s(\pmb{a})=(\pmb{a}\cup\pmb{b})$ and
  $s(\pmb{b})=(\pmb{a}\cup\pmb{b})$.  Every word in $L$ has even length, so
  every word in $\overline{s}(L)$ also has even length.  Thus
  $\overline{s}(L)$ is contained in the set of even-length words.

  Conversely, let $y$ be any even-length word, say $|y|=2n$.  Choose
  $x=\pmb{a}^n\pmb{b}^n\in L$.  Since each letter of $x$ may be replaced by
  either $\pmb{a}$ or $\pmb{b}$, the word $y$ belongs to $\overline{s}(x)$.
  Hence every even-length word lies in $\overline{s}(L)$.

  Therefore $\overline{s}(L)$ is exactly the set of all even-length words over
  $\{\pmb{a},\pmb{b}\}$, which is regular.  So the image of a nonregular
  language under a substitution need not remain in $\mathfrak{N}_\Gamma$.
\end{enumerate}

%% 6.45
\item \textbf{Proof of Lemma~6.3.}

Lemma~6.3 is the version of Theorem~6.2 that eliminates the $m$th unknown
instead of necessarily the last unknown.

\begin{proof}
\emph{Proof.}  Renumber the unknowns so that $X_m$ becomes the last unknown
$X_n$, without changing the actual system except for notation.  After this
relabeling, Theorem~6.2 applies directly.  It gives the elimination formulas
\[
\hat{E}_i=E_i\cup(A_{im}A_{mm}^*E_m)
\qquad\text{and}\qquad
\hat{A}_{ij}=A_{ij}\cup(A_{im}A_{mm}^*A_{mj}),
\]
for every index other than $m$, together with the back-substitution formula
\[
X_m=A_{mm}^*\bigl(E_m\cup A_{m1}X_1\cup\cdots\cup A_{m(m-1)}X_{m-1}
\cup A_{m(m+1)}X_{m+1}\cup\cdots\cup A_{mn}X_n\bigr).
\]
Renaming the variables back proves Lemma~6.3. \qed
\end{proof}

%% 6.46
\item $\Xi=\{x\in\{\pmb{a},\pmb{b}\}^*\mid x\text{ contains at least two consecutive }\pmb{b}\text{s}\wedge x\text{ contains no two consecutive }\pmb{a}\text{s}\}$:

Words over $\{\pmb{a},\pmb{b}\}$ avoiding $\pmb{aa}$ are exactly
$(\pmb{b}\cup\pmb{ab})^*(\pmb{a}\cup\epsilon)$.  Requiring a $\pmb{bb}$
substring gives:
\[
  (\pmb{b}\cup\pmb{ab})^*(\pmb{a}\cup\epsilon)\pmb{bb}(\pmb{b}\cup\pmb{ab})^*(\pmb{a}\cup\epsilon)
\]

%% 6.47
\item \begin{enumerate}[(a)]
  \item Every $\pmb{b}$ eventually followed by $\pmb{c}$ (not necessarily immediately):
  A word is in the language iff either it has no $\pmb{b}$ at all after its
  last $\pmb{c}$, or it has no $\pmb{c}$ and then must contain only $\pmb{a}$s.
  Thus the suffix after the last $\pmb{c}$ must lie in $\pmb{a}^*$:
  \[
    \pmb{a}^* \cup (\pmb{a}\cup\pmb{b}\cup\pmb{c})^*\pmb{c}\pmb{a}^*
  \]
  Any $\pmb{b}$ before the last $\pmb{c}$ is discharged by that final
  $\pmb{c}$, and after the last $\pmb{c}$ only $\pmb{a}$s may remain.

  \item Every $\pmb{b}$ immediately followed by $\pmb{c}$: (same as Exercise~6.11)
  \[
    (\pmb{a}\cup\pmb{c}\cup\pmb{bc})^*
  \]
\end{enumerate}

%% 6.48
\item Over $\Sigma=\{\pmb{a},\pmb{b}\}$:

\begin{enumerate}[(a)]
  \item No consecutive $\pmb{a}$s and no consecutive $\pmb{b}$s (strictly alternating):
  \[
    (\pmb{ab})^*(\pmb{a}\cup\epsilon) \cup (\pmb{ba})^*(\pmb{b}\cup\epsilon)
  \]
  (Alternating starting with $\pmb{a}$ or $\pmb{b}$, or empty word.)

  \item Words beginning and ending with different letters:
  \[
    \pmb{a}(\pmb{a}\cup\pmb{b})^*\pmb{b} \cup \pmb{b}(\pmb{a}\cup\pmb{b})^*\pmb{a}
  \]

  \item Words for which the last two letters match ($\pmb{aa}$ or $\pmb{bb}$ at the end):
  \[
    (\pmb{a}\cup\pmb{b})^*\pmb{aa} \cup (\pmb{a}\cup\pmb{b})^*\pmb{bb}
  \]

  \item Words for which the first two letters match ($\pmb{aa}$ or $\pmb{bb}$ at the start):
  \[
    \pmb{aa}(\pmb{a}\cup\pmb{b})^* \cup \pmb{bb}(\pmb{a}\cup\pmb{b})^*
  \]

  \item Words for which the first and last letters match:
  \[
    \epsilon \cup \pmb{a} \cup \pmb{b} \cup \pmb{a}(\pmb{a}\cup\pmb{b})^*\pmb{a} \cup \pmb{b}(\pmb{a}\cup\pmb{b})^*\pmb{b}
  \]
\end{enumerate}

%% 6.49
\item \textbf{The language of all valid regular expressions over $\{\pmb{a},\pmb{b}\}$ is not FAD.}

\begin{proof}
Let $V$ be the language of all valid regular expressions over
$\{\pmb{a},\pmb{b}\}$.  Intersect $V$ with the regular language
\[
R=(\pmb{(})^*\pmb{a}(\pmb{)})^*.
\]
The only valid regular expressions in $R$ are the strings
\[
\pmb{(}^n\pmb{a}\pmb{)}^n \qquad (n\ge 0),
\]
because the parentheses must match and simply enclose the one-letter expression
$\pmb{a}$.  Therefore
\[
V\cap R=\{\pmb{(}^n\pmb{a}\pmb{)}^n\mid n\ge 0\},
\]
which is not FAD by the pumping lemma.  But the intersection of two FAD
languages would be FAD.  Hence $V$ is not FAD.
\end{proof}

%% 6.50
\item Apply Theorem~6.3 (state equations) and Theorem~6.2 (Arden's lemma) to each NDFA in Figure~6.14 to produce a regular expression.

\begin{enumerate}[(a)]
  \item $(\pmb{a}\cup\pmb{c})^*$
  \item $(\pmb{b}\cup\pmb{a}\pmb{c}^*\pmb{b})^*$
  \item $(\pmb{a}\cup(\pmb{a}\cup\pmb{c})\pmb{b})^*$
  \item $\pmb{a}(\pmb{bb})^*\pmb{a}$
  \item $(\pmb{a}\cup\pmb{aa}\cup\pmb{ab})^*\pmb{a}$
  \item $\pmb{a}\pmb{a}^*\pmb{b}\pmb{c}(\pmb{a}^*\pmb{b}\pmb{c})^*$
  \item $\bigl(\pmb{a}(\pmb{bb}\cup\pmb{cc})^*\pmb{a}\bigr)^*$
\end{enumerate}

%% 6.51
\item \textbf{Complete details of Theorem~6.6.}

Theorem~6.6 proves closure of \RSIG\ under substitution.  The text supplies the
basis step; the missing details are the three operator cases.

\begin{proof}
Assume inductively that whenever $R_1$ and $R_2$ have at most $m$ operators, the
expressions obtained by replacing each letter $\pmb{a}$ with $s(\pmb{a})$
represent $\overline{s}(R_1)$ and $\overline{s}(R_2)$.

If $R=R_1\cdot R_2$, then by Definition~6.5,
\[
\overline{s}(R_1\cdot R_2)=\overline{s}(R_1)\cdot\overline{s}(R_2).
\]
So the substituted expression $R_1'R_2'$ represents $\overline{s}(R)$.

If $R=R_1\cup R_2$, then
\[
\overline{s}(R_1\cup R_2)=\overline{s}(R_1)\cup\overline{s}(R_2),
\]
so the substituted expression $R_1'\cup R_2'$ represents $\overline{s}(R)$.

If $R=(R_1)^*$, then
\[
\overline{s}((R_1)^*)=(\overline{s}(R_1))^*,
\]
so $(R_1')^*$ represents $\overline{s}(R)$.

Thus substitution preserves each of the three operators used to build regular
expressions.  By induction on the number of operators in $R$, the expression
obtained by substituting $s(\pmb{a})$ for each letter $\pmb{a}$ represents
exactly $\overline{s}(R)$.  Therefore \RSIG\ is closed under substitution. \qed
\end{proof}

%% 6.52
\item \textbf{Prove Lemma~6.4.}

\begin{proof}
Let
\[
A=\langle\Sigma,\{s_1,s_2,\ldots,s_n\},S_0,\delta,F\rangle
\]
be an NDFA, and for each $i$ let
\[
Y_i=I(A,s_i)=\{x\mid \DBAR(S_0,x)\ni s_i\}.
\]
Define
\[
I_i=
\begin{cases}
\epsilon & \text{if } s_i\in S_0,\\
\emptyset & \text{if } s_i\not in S_0,
\end{cases}
\]
and
\[
B_{ij}=\bigcup_{\pmb{a}\in\Sigma \wedge  s_i\in\delta(s_j,\pmb{a})}\pmb{a}.
\]
We show that
\[
Y_k=I_k\cup Y_1B_{k1}\cup Y_2B_{k2}\cup\cdots\cup Y_nB_{kn}.
\]

If $x\in Y_k$, then either $x=\lambda$ and $s_k\in S_0$, so $x\in I_k$, or
$x=za$ where $z$ takes some start state to some $s_j$ and $a$ takes $s_j$ to
$s_k$.  Then $z\in Y_j$ and $a\in B_{kj}$, so $x\in Y_jB_{kj}$.  Hence
$Y_k\subseteq I_k\cup Y_1B_{k1}\cup\cdots\cup Y_nB_{kn}$.

Conversely, if $x\in I_k$, then $x=\lambda$ and $s_k\in S_0$, so $x\in Y_k$.
If $x\in Y_jB_{kj}$, then $x=za$ for some $z\in Y_j$ and some letter
$a\in B_{kj}$.  Thus $z$ reaches $s_j$ from a start state, and then $a$ takes
$s_j$ to $s_k$, so $x\in Y_k$.  Therefore
$I_k\cup Y_1B_{k1}\cup\cdots\cup Y_nB_{kn}\subseteq Y_k$ as well.

So the required equations hold for every $k$. \qed
\end{proof}

%% 6.53
\item \textbf{Theorems~5.2, 5.4, 5.5 as corollaries of Theorem~6.6.}

By Theorem~6.5, every FAD language has a regular expression.  Let
$L_1=L(R_1)$ and $L_2=L(R_2)$.

For union, work over the alphabet $\{\pmb{a},\pmb{b}\}$ with the regular
expression
\[
Q=\pmb{a}\cup\pmb{b}.
\]
Define a substitution by
\[
s(\pmb{a})=R_1,\qquad s(\pmb{b})=R_2.
\]
Then the substituted expression is
\[
\overline{s}(Q)=R_1\cup R_2,
\]
whose language is $L_1\cup L_2$.  So $L_1\cup L_2$ is regular by Theorem~6.6.

For concatenation, use the regular expression
\[
Q=\pmb{ab}
\]
with the same substitution.  Then
\[
\overline{s}(Q)=R_1R_2,
\]
whose language is $L_1\cdot L_2$.  So $L_1\cdot L_2$ is regular.

For Kleene closure, work over the alphabet $\{\pmb{a}\}$ and use the regular
expression
\[
Q=\pmb{a}^*.
\]
Define
\[
s(\pmb{a})=R_1.
\]
Then
\[
\overline{s}(Q)=(R_1)^*,
\]
whose language is $L_1^*$.  So $L_1^*$ is regular.

Thus Theorems~5.2, 5.4, and 5.5 follow as corollaries of Theorem~6.6.

%% 6.54
\item \textbf{Class $\pmb{F}$ generated by the five rules:}

The five rules generate exactly the class of \emph{finite} languages over $\{\pmb{a},\pmb{b}\}$:
\begin{enumerate}[(i)]
  \item Singletons $\{\pmb{a}\}$ and $\{\pmb{b}\}$ are in $\pmb{F}$.
  \item $\emptyset\in\pmb{F}$.
  \item $\{\lambda\}\in\pmb{F}$.
  \item Concatenation: $F_1\cdot F_2$ is finite if $F_1$ and $F_2$ are.
  \item Union: $F_1\cup F_2$ is finite if $F_1$ and $F_2$ are.
\end{enumerate}
No Kleene closure is allowed, so $\pmb{F}$ consists only of finite languages.  Every finite language over $\{\pmb{a},\pmb{b}\}$ is in $\pmb{F}$ (can be expressed as a finite union of finite concatenations of singletons), and $\pmb{F}\subseteq F_{\{\pmb{a},\pmb{b}\}}$ (the finite languages).  Thus $\pmb{F}=F_{\{\pmb{a},\pmb{b}\}}$, the class of all finite languages over $\{\pmb{a},\pmb{b}\}$.

\end{enumerate}
